\begin{proof}
Demostrado este lema, la demostración de (6.14.1) se desarrolla en varias
etapas.
I) \emph{Reducción al caso en que $Y=\operatorname{Spec}(A)$,
$Y'=\operatorname{Spec}(A')$, $X=\operatorname{Spec}(B)$, donde $A$ y $A'$
son anillos locales noetherianos, $A$ es un dominio entero y $B$ es la
clausura integral de $A$.} --- Hay que demostrar que, para todo
$x'\in X'$, el anillo local $\mathcal O_{X',x'}$ es un dominio entero
integralmente cerrado; sean $x$, $y$, $y'$ las imágenes canónicas de
$x'$ en $X$, $Y$, $Y'$, respectivamente; si se pone
$A=\mathcal O_{Y,y}$, $A'=\mathcal O_{Y',y'}$,
$B=\mathcal O_{X,x}$, los anillos $\mathcal O_{X',x'}$ para $x$, $y$,
$y'$ fijos son los anillos locales en los ideales primos del anillo
$B'=B\otimes_A A'$ (\textbf I, 3.6.5), por lo que hay que demostrar que $B'$
es un anillo \emph{normal}. Observemos, por otra parte, que
$\operatorname{Spec}(A')\to\operatorname{Spec}(A)$ es un morfismo normal
por definición (6.8.1 e \textbf{I}, 3.6.5); la cuestión se reduce, pues, al caso en que
$Y$ e $Y'$ son esquemas locales y $X$ es el espectro de un anillo local
\emph{íntegro e integralmente cerrado} $B$. Denotemos por
$(B_\alpha)$ la familia de las \emph{clausuras integrales} de las
sub-$A$-álgebras \emph{de tipo finito} de $B$; está claro que $B$ es la unión
de la familia filtrante creciente de los $B_\alpha$. Como el funtor
$\varinjlim$ conmuta con el producto tensorial, $B'$ es, por tanto,
isomorfo a $\varinjlim B'_\alpha$, donde se ha puesto
$B'_\alpha=B_\alpha\otimes_A A'$. Para demostrar que $B'$ es normal,
bastará, en virtud de (5.13.6), demostrar que los anillos $B'_\alpha$
son \emph{normales} y que, para $\alpha\leqslant\beta$, toda componente
irreducible de $\operatorname{Spec}(B'_\beta)$ \emph{domina} una componente
irreducible de $\operatorname{Spec}(B'_\alpha)$. Pero esta última
propiedad se deduce de la hipótesis de que $A'$ es un $A$-módulo plano
y de (2.3.7, (ii)), puesto que $B_\alpha$ y $B_\beta$ son dominios
enteros y $B_\alpha\subset B_\beta$.
